\documentclass[11pt, a4paper]{article}

% ------------------------------------------------------------
% Packages
% ------------------------------------------------------------
\usepackage[a4paper, top=2.6cm, bottom=2.6cm, left=2.5cm, right=2.5cm]{geometry}
\usepackage[utf8]{inputenc}
\usepackage[T1]{fontenc}
\usepackage[english]{babel}
\usepackage[sfdefault]{noto}          % Noto Sans (pdfLaTeX)
\renewcommand{\familydefault}{\sfdefault}

\usepackage{amsmath, amssymb, bm}
\usepackage{booktabs, array, tabularx}
\usepackage{enumitem}
\usepackage{xcolor}
\usepackage{microtype}
\usepackage{titlesec}
\usepackage{fancyhdr}
\usepackage{setspace}
\usepackage[most]{tcolorbox}
\usepackage[colorlinks=true, linkcolor=accent, urlcolor=accent, citecolor=accent,
            pdftitle={Multi-Scene VGGT-Omega: Project Plan},
            pdfauthor={}]{hyperref}

% ------------------------------------------------------------
% Colours and layout
% ------------------------------------------------------------
\definecolor{accent}{HTML}{1F4E79}
\definecolor{softgray}{HTML}{F3F5F8}
\definecolor{rulegray}{HTML}{C9D1DB}

\setstretch{1.12}
\setlength{\parindent}{0pt}
\setlength{\parskip}{0.55em}
\setlist{itemsep=0.25em, topsep=0.3em, parsep=0pt}
\setlist[itemize,1]{label=\textcolor{accent}{\textbullet}}
\setlist[enumerate,1]{label=\textcolor{accent}{\textbf{\arabic*.}}}

% Section headings
\titleformat{\section}
  {\Large\bfseries\color{accent}}{\thesection}{0.7em}{}
  [\vspace{-0.4em}{\color{rulegray}\titlerule[0.8pt]}]
\titleformat{\subsection}
  {\large\bfseries\color{accent!85!black}}{\thesubsection}{0.6em}{}
\titlespacing*{\section}{0pt}{1.6em}{0.8em}
\titlespacing*{\subsection}{0pt}{1.2em}{0.4em}

% Header / footer
\pagestyle{fancy}
\fancyhf{}
\renewcommand{\headrulewidth}{0.4pt}
\renewcommand{\headrule}{\hbox to\headwidth{\color{rulegray}\leaders\hrule height \headrulewidth\hfill}}
\fancyhead[L]{\small\color{gray} Multi-Scene VGGT-$\Omega$}
\fancyhead[R]{\small\color{gray} Project Plan}
\fancyfoot[C]{\small\color{gray} \thepage}

% Info box
\newtcolorbox{keybox}[1][]{
  enhanced, breakable, colback=softgray, colframe=accent,
  boxrule=0pt, leftrule=3pt, arc=0pt, left=10pt, right=10pt, top=6pt, bottom=6pt,
  fonttitle=\bfseries\color{accent}, coltitle=accent, title={#1},
  attach boxed title to top left={yshift=-2mm, xshift=0mm},
  boxed title style={colback=softgray, boxrule=0pt, left=0pt, top=0pt, bottom=0pt},
  before skip=10pt, after skip=10pt
}

% Macros
\newcommand{\zcam}{\bm{z}_{\mathrm{cam}}}
\newcommand{\zscene}{\bm{z}_{\mathrm{scene}}}
\newcommand{\zdino}{\bm{z}_{\mathrm{dino}}}
\newcommand{\Lall}{\mathcal{L}}
\newcommand{\VGGT}{VGGT-$\Omega$}

% ------------------------------------------------------------
% Title
% ------------------------------------------------------------
\title{%
  {\Huge\bfseries\color{accent} Multi-Scene \VGGT}\\[0.6em]
  {\Large\color{gray} Detailed Project Plan and Architecture Strategy}%
}
\author{}
\date{}

\begin{document}

\maketitle
\thispagestyle{fancy}
\vspace{-2.5em}

\begin{keybox}[Summary]
This document describes the project plan, architecture, and training strategy of the
\textbf{Multi-Scene \VGGT} network. The model appends \emph{scene} and \emph{camera}
register tokens to DINO features, applies alternating attention across frames, groups
the input images into scenes with an internal affinity predictor, and jointly solves
for 3D geometry and camera poses of every discovered scene.
\end{keybox}

\tableofcontents
\vspace{1em}

% ============================================================
\section{Architecture Overview}
% ============================================================
The architecture processes features from a foundation model (DINO) and augments them with
specialized register tokens that model camera and scene information.

\subsection{Tokenization and Feature Extraction}
For each input image $i$, spatial features $\zdino$ are extracted with DINO and
concatenated with two sets of learnable register tokens:
\begin{itemize}
  \item \textbf{Camera token} $\zcam^{(i)} \in \mathbb{R}^{1 \times C}$: a single token that
        registers camera-specific information.
  \item \textbf{Scene token} $\zscene^{(i)} \in \mathbb{R}^{16 \times C}$: a group of 16 tokens
        that captures global scene context.
\end{itemize}
The final representation of image $i$ is
\begin{equation}
  \bm{z}_i = \big[\, \zdino^{(i)} \,;\, \zcam^{(i)} \,;\, \zscene^{(i)} \,\big]
  \in \mathbb{R}^{(HW + 17) \times C},
\end{equation}
where $H \times W$ is the spatial resolution of the feature map and $C$ is the channel
dimension.

\subsection{Alternating Attention}
The token sequence is processed by the \VGGT{} transformer, which consists of
\textbf{24 blocks}. Each block alternates between two attention types:
\begin{enumerate}
  \item \textbf{Frame attention:} self-attention computed \emph{within} each frame
        (intra-frame).
  \item \textbf{Global attention:} attention computed across the tokens of \emph{all}
        frames in the batch (inter-frame).
\end{enumerate}

% ============================================================
\section{Scene Attention and Affinity Prediction}
% ============================================================
Frames that belong to the \emph{same} scene should exchange information, while unrelated
frames must remain isolated. This is achieved by a dedicated \textbf{scene attention}
step, applied after the 24 alternating-attention blocks (referred to as block 25).

\subsection{Standard Attention over Scene Tokens}
The scene tokens of all frames are gathered. For a batch of $N$ images the total number of
scene tokens is $M = 16N$. The query, key, and value matrices are
\begin{equation}
  \bm{Q}, \bm{K}, \bm{V} \in \mathbb{R}^{M \times d_k},
\end{equation}
and the base attention score matrix is
\begin{equation}
  \bm{S} = \frac{\bm{Q}\bm{K}^{\top}}{\sqrt{d_k}} \in \mathbb{R}^{M \times M}.
\end{equation}

\subsection{Affinity Matrix and Masking}
To block attention between tokens of \emph{different} scenes, a pairwise affinity matrix
$\hat{\bm{A}} \in [0,1]^{N \times N}$ is learned in four steps:

\begin{enumerate}
  \item \textbf{Descriptor extraction.} Average the 16 scene tokens of each image to obtain a
        scene descriptor
        \begin{equation}
          \bm{h}_i = \frac{1}{16}\sum_{k=1}^{16} \zscene^{(i)}[k] \in \mathbb{R}^{C}.
        \end{equation}
  \item \textbf{Affinity prediction.} Pass each pair of descriptors $(\bm{h}_i, \bm{h}_j)$
        through a linear layer followed by a sigmoid to obtain $\hat{A}_{ij}$:
        \begin{itemize}
          \item $\hat{A}_{ij} \approx 1$: images $i$ and $j$ belong to the \textbf{same} scene.
          \item $\hat{A}_{ij} \approx 0$: images $i$ and $j$ belong to \textbf{different} scenes.
        \end{itemize}
  \item \textbf{Thresholding.} Convert $\hat{\bm{A}}$ into an additive mask
        $\bm{B} \in \{0, -\infty\}^{N \times N}$ with threshold $\tau$:
        \begin{equation}
          B_{ij} =
          \begin{cases}
            0       & \text{if } \hat{A}_{ij} \ge \tau \quad (\text{same scene}),\\
            -\infty & \text{otherwise} \quad (\text{different scene}).
          \end{cases}
        \end{equation}
        The mask is expanded to token level as $\bm{B} \otimes \bm{1}_{16 \times 16}
        \in \mathbb{R}^{M \times M}$, so that every image pair controls a $16 \times 16$
        block of attention scores.
  \item \textbf{Masked attention.} Add the mask to the scores before the softmax:
        \begin{equation}
          \bm{W} = \mathrm{Softmax}\!\left(\bm{S} + \bm{B} \otimes \bm{1}_{16 \times 16}\right).
        \end{equation}
        Entries set to $-\infty$ receive an attention weight of exactly $0$, which strictly
        prevents cross-scene contamination.
\end{enumerate}

% ============================================================
\section{Scene Grouping and Forward Output}
% ============================================================
After thresholding, the batch of independent images is clustered into distinct scene
groups, for example:
\begin{itemize}
  \item \textbf{Group A:} images $\{1, 2, 3\}$
  \item \textbf{Group B:} images $\{4, 5, 6\}$
\end{itemize}
For each group an \textbf{origin (reference) image} is assigned. It serves as the anchor
for reconstructing the 3D scene and for computing relative camera poses.

The forward pass returns its output hierarchically, one entry per scene:
\begin{tcolorbox}[enhanced, colback=softgray, colframe=rulegray, boxrule=0.5pt, arc=2pt,
                  left=8pt, right=8pt, top=4pt, bottom=4pt]
\begin{verbatim}
output = [
    {
        "images":       [1, 2, 3],
        "pointcloud":   ...,   # 3D points of scene A
        "camera_poses": ...    # poses relative to the origin image
    },
    {
        "images":       [4, 5, 6],
        ...
    }
]
\end{verbatim}
\end{tcolorbox}

% ============================================================
\section{Loss Functions}
% ============================================================
The training objective is a weighted sum of three terms that jointly optimize scene
clustering, geometric reconstruction, and camera localization.

\begin{table}[h]
  \centering
  \small
  \renewcommand{\arraystretch}{1.35}
  \begin{tabularx}{\textwidth}{@{}l l X@{}}
    \toprule
    \textbf{Term} & \textbf{Symbol} & \textbf{Purpose} \\
    \midrule
    Affinity loss & $\Lall_{\text{affinity}}$ & Supervises the affinity head so that scene groupings are identified correctly. \\
    Geometry loss & $\Lall_{\text{geom}}$ & 3D point error of the reconstruction of each separated scene. \\
    Pose loss     & $\Lall_{\text{pose}}$ & Rotation and translation error w.r.t.\ each scene's anchor image. \\
    \bottomrule
  \end{tabularx}
\end{table}

\subsection{Affinity Loss}
The predicted matrix $\hat{\bm{A}}$ is compared with the ground-truth adjacency matrix
$\bm{A}$ (where $A_{ij}=1$ if images $i$ and $j$ come from the same scene):
\begin{equation}
  \Lall_{\text{affinity}} = \mathrm{BCE}\big(\hat{\bm{A}}, \bm{A}\big).
\end{equation}

\subsection{Geometry Loss}
$\Lall_{\text{geom}}$ is a 3D point loss evaluated on the reconstructed geometry of each
separated scene. It is computed as a per-pixel 3D distance between predicted and ground-truth
point clouds (or depth maps).

\subsection{Pose Loss}
$\Lall_{\text{pose}}$ is evaluated independently for each scene with respect to its anchor
image and combines rotational and translational errors:
\begin{equation}
  \Lall_{\text{pose}} = \Lall_{\text{rot}} + \lambda_t \, \Lall_{\text{trans}}.
\end{equation}

\subsection{Total Loss}
The three tasks are balanced with hyperparameters $\lambda_1, \lambda_2, \lambda_3$:
\begin{equation}
  \Lall_{\text{total}} =
  \lambda_1 \, \Lall_{\text{affinity}} +
  \lambda_2 \, \Lall_{\text{geom}} +
  \lambda_3 \, \Lall_{\text{pose}}.
\end{equation}

% ============================================================
\section{Training Strategy and Dataset}
% ============================================================

\subsection{Mixed-Scene Batch Generation}
To train the affinity predictor, every batch must contain multiple scenes. A
\textbf{mixed-scene batch generator} builds batches of size
\begin{equation}
  N = S \times K,
\end{equation}
where
\begin{itemize}
  \item $N$ is the total number of frames in the batch,
  \item $S$ is the number of distinct scenes sampled, and
  \item $K$ is the number of frames per scene, with the strict condition $K \ge 3$.
\end{itemize}

\subsection{Datasets}
The model is trained on a diverse mix of indoor and outdoor multi-view datasets:

\begin{table}[h]
  \centering
  \small
  \renewcommand{\arraystretch}{1.3}
  \begin{tabular}{@{}l l@{}}
    \toprule
    \textbf{Dataset} & \textbf{Setting} \\
    \midrule
    ScanNet / ScanNet++ & Indoor \\
    Replica             & Indoor \\
    MegaDepth           & Outdoor \\
    ETH3D               & Indoor / outdoor \\
    \bottomrule
  \end{tabular}
\end{table}

\subsection{Two-Stage Fine-Tuning}
To stabilize the multi-task optimization, training is split into two stages:
\begin{enumerate}
  \item \textbf{Stage 1: Affinity initialization.} Freeze the transformer backbone and
        train only the affinity head with $\Lall_{\text{affinity}}$. The network must first
        separate scenes reliably before geometry and pose are estimated.
  \item \textbf{Stage 2: End-to-end fine-tuning.} Unfreeze all parameters and train end to end
        with the full $\Lall_{\text{total}}$, using a significantly smaller learning rate.
\end{enumerate}

\end{document}
